\section{DynamicMCPBench}
\label{sec:benchmark}

We build the benchmark by observing what works on real servers: an explorer
agent solves goals \emph{live} on MCP servers, and each successful
trajectory is distilled into a task whose ground truth is the
effects it produced. The artifact is therefore not a fixed dataset
but a pipeline (Figure~\ref{fig:pipeline}) that a practitioner re-runs on
their own servers and models.

\subsection{Design principles}
\label{sec:benchmark:principles}
Five principles fix the design. (1)~\emph{The trace is the primitive}: tasks
come from recorded successful trajectories, so every required tool was
actually used and a spurious ``unnecessary'' tool cannot arise. (2)~\emph{We
never grade the final answer}: scoring checks effects, not whether a
free-form reply matches a reference string, a match that is unstable the
moment the answer depends on live data. (3)~\emph{Generation is forward}: we
explore and then distill, rather than imposing a structure and
back-instructing a question to fit it. (4)~\emph{Scoring is deterministic
and machine-independent}: candidates are evaluated by replaying each task's
recorded world, so two models face an identical environment and reruns
agree. (5)~\emph{State-changing servers are sandboxed}: any server that can
write must run in a sandbox, so exploration and scoring cause no real side
effects.

\paragraph{Why ``dynamic''.} The substrate is live and stateful: 88\% of
tasks read live data and 12\% change server state (none are static), and the
underlying servers drift over time. We freeze each task's world through
deterministic replay for fair scoring, while a refresh step re-runs the
reference trajectories against the live servers and flags any task whose
effects can no longer be reproduced.

\paragraph{Safety.} Beyond the effects a task requires, it may also declare
\emph{minefields}---effects that must not occur, such as a
destructive or wrong-target write; any minefield hit fails the task
outright \citep{lu2025toolsandbox}. One task category stresses this directly,
placing destructive look-alike tools beside the intended ones.

\subsection{Substrate and corpus}
\label{sec:benchmark:substrate}
The released corpus contains 1{,}845 tasks distilled from 2{,}051 reference
trajectories over 121 live MCP servers, and it is authored by a diverse pool
of frontier model families (the largest contributing 18\% of the tasks) so
that no single generator shapes it. Tasks span 15 \emph{categories}, each
targeting a distinct tool-use challenge of the generated question, from a
neutral baseline through semantic near-misses, cross-server and multi-step
structure, ordering and recovery demands, traps, and under-specification;
Appendix~\ref{app:categories} defines all fifteen, and
Appendix~\ref{app:corpus} reports the corpus composition. Tasks are further stratified by
\emph{dynamism} (live-read vs.\ state-changing), by \emph{length}, the depth
of the tool-dependency chain, which ranges from~1 to~77 with a mean
of~5.2, and by \emph{scope}, intra- vs.\ cross-server (29\% of tasks span
more than one server). For the study we evaluate on a balanced slice of 50
tasks per category, 750 in total.

\subsection{Forward generation and distillation}
\label{sec:benchmark:gen}
A goal generator turns each server's tool surface into realistic user goals,
with persona-varied phrasing for diversity; an explorer agent then pursues
one goal live, and a successful run is handed to a distiller that emits a
\emph{task specification}. A specification has a fuzzy natural-language
prompt (tool names removed); one or more \emph{checkpoints} of two kinds---a
\emph{tool-effect} checkpoint (some tool from an \emph{equivalence set} of
interchangeable tools must be called, optionally meeting an argument
constraint) and a \emph{value-produced} checkpoint (a demanded value must
appear in a tool result); optional \emph{minefields}; and a \emph{partial
order} that constrains two effects only where one genuinely depends on the
other, leaving parallel steps unordered. Each specification also records a
complexity profile (chain depth, cross-server, branching) used for the
stratification above; Table~\ref{tab:worked-example} works one task through end to end, from
recorded trajectory to distilled checkpoints; Appendix~\ref{app:example}
shows the full TaskSpec rendering. The choices that shape the corpus are deliberate:
generation is forward and trace-grounded so tasks are provably achievable;
categories are balanced at 50 tasks each so no category dominates; seed
difficulty is scaled to span short and long chains; and authorship is spread
across model families. The distiller never introduces a checkpoint the trace
does not justify. The full generation and evaluation configuration appears in
Appendix~\ref{app:config}, and all prompts in Appendix~\ref{app:prompts}.


\begin{table*}[t]
\centering
\small
\setlength{\tabcolsep}{4pt}
\renewcommand{\arraystretch}{1.05}
\caption{A compact worked example. A successful reference trace is distilled
into path-agnostic effect checkpoints. A candidate passes by reproducing these
effects with any equivalent tool, not by matching the final answer text.}
\label{tab:worked-example}
\begin{tabular}{p{0.18\linewidth}p{0.76\linewidth}}
\toprule
\textbf{Stage} & \textbf{Example} \\
\midrule

User goal &
Compare the financial health and recent performance of Apple (AAPL),
Microsoft (MSFT), and Google (GOOGL): latest ticker information, most recent
quarterly earnings, and one year of price history. \\

Reference trace &
The explorer calls \texttt{get\_tickers\_info} for the three symbols,
\texttt{get\_earnings} for each company, \texttt{download} for one-year price
history, and \texttt{get\_financials} for yearly balance-sheet and income
statements. \\

Distilled effects &
The TaskSpec requires ticker information for all three symbols, quarterly
earnings, one-year price history, yearly balance-sheet financials, yearly
income-statement financials, and a final message mentioning the companies and
the relevant financial evidence. \\

Equivalence &
The price-history checkpoint accepts either \texttt{download} or
\texttt{get\_price\_history}, so an agent need not repeat the reference path
if it obtains the same effect through an equivalent tool. \\

Minefields / order &
Minefields: none. Partial order: none, because the required evidence can be
collected independently. \\

Scoring &
A candidate passes if all required effects are observed under deterministic
replay in all three attempts. It fails if an effect is missing, even if the
final natural-language answer looks plausible. \\

\bottomrule
\end{tabular}
\end{table*}

\subsection{Effect-based scoring}
\label{sec:benchmark:eval}

Given a candidate trajectory, a deterministic first tier (Tier-1) checks each
checkpoint: that some tool in a tool-effect's equivalence set was called with
arguments satisfying its constraint; that a value-produced checkpoint appears;
that no minefield was hit; and that the partial order holds. Because
checkpoints carry equivalence sets, any path achieving the required effects
passes. On the released corpus, 16\% of effect checkpoints admit two or more
interchangeable tools (up to twelve), so multiple valid trajectories are
accepted rather than a single gold path (Appendix~\ref{app:eqset}). A second
tier (Tier-2) is an LLM judge that may upgrade only a \emph{failed}
tool-effect when a different tool is genuinely effect-equivalent; it never
reads or grades the final answer. The headline leaderboard uses Tier-1 alone
(Appendix~\ref{app:config}).

Candidates are compared under deterministic replay. Each task is attempted
three times and counted as solved only if all three attempts pass
(pass\textasciicircum{}3; \citealp{yao2024tau}). To keep the setting controlled and identical
across models, each candidate receives the required tools plus fixed
distractors, half of them same-name tools on other servers. The action budget
covers the longest chains, and the main configuration presents tool
descriptions as written and exposes all tools directly, so the headline
measures the agent rather than a description rewriter or retriever.


\subsection{Applying the framework to your own servers}
\label{sec:benchmark:apply}

Because every stage (collecting servers, generating goals, exploring,
distilling, and scoring) is automated, the same pipeline runs end to end on a
practitioner's own servers and chosen models. In an industrial deployment, a
team can point DynamicMCPBench at its private MCP server fleet, generate tasks
that reflect its own workflows, and obtain a deployment-specific benchmark,
leaderboard, and failure breakdown. That is the mode behind the study in \S\ref{sec:results}.

The public benchmark then serves as a calibrated reference point. A company can
compare candidate agents not only by an aggregate score, but by the task
profile that matters for its own stack: short versus long tool chains,
single-server versus cross-server composition, recovery requirements, or
look-alike-tool confusion. Thus the framework supports model selection and
regression testing for private MCP deployments rather than only a static public
leaderboard.

